\section{Residual Networks}
\label{sec:resnet}

The convolutional networks of Section~\ref{sec:cnn} became better with every
added layer, from the eight layers of AlexNet to the nineteen of VGG. It is
tempting to continue: if $19$ layers beat $8$, then $50$ or $200$ layers should
be better still. In practice, simply stacking more layers soon stops working.
Signals shrink or explode on their way through the network, gradients become
erratic, and beyond a certain depth even the \emph{training} error goes up.

Three techniques together make very deep networks trainable, and all modern
image networks use them. \emph{Careful parameter initialization} keeps the size
of activations and gradients constant from layer to layer at the start of
training. \emph{Residual connections} add the input of each block to its output,
which gives the gradient a direct path through the network. \emph{Batch
normalization} rescales activations during training so that they stay in a
well-behaved range. This section explains each technique, starting with the
problem it solves, and then presents the architectures built from them: ResNet,
DenseNet, U-Net, and stacked hourglass networks.

\subsection{Initialization and the Variance of Activations}

Training starts from randomly chosen parameters. A common choice sets all biases
to zero and draws every weight from a normal distribution with mean zero and
variance $\sigma_\Omega^2$. The value of $\sigma_\Omega^2$ looks like a minor
detail, but in a deep network it decides whether training can start at all.
Every layer multiplies the signal by a weight matrix. If the weights are small,
the activations of the next layer are likely to be smaller than those of the
current layer; after fifty layers they have practically vanished. If the weights
are large, the activations grow from layer to layer and eventually overflow.
The backward pass also multiplies by the (transposed) weight matrices, so the
same happens to the gradients. This is the \emph{vanishing gradient} and the
\emph{exploding gradient} problem: the early layers receive gradients that are
either far too small to change anything or so large that every update destroys
what was learned.

\paragraph{How the variance propagates.}
The right value of $\sigma_\Omega^2$ follows from tracking the variance of the
signal through one layer. Let $f_j$ be the pre-activations of layer $k$, with
mean zero and variance $\sigma_f^2$. The layer applies ReLU,
$h_j = \max(0, f_j)$, and the next layer computes the pre-activations
\begin{equation}
  f'_i = \beta_i + \sum_{j=1}^{D_h} \Omega_{ij} h_j ,
\end{equation}
where $D_h$ is the number of units in layer $k$. With zero biases and weights
that are independent of the activations, the variance of $f'_i$ is the sum of
the variances of the $D_h$ products. Each product has variance
$\sigma_\Omega^2\,\mathbb{E}[h_j^2]$. Because $f_j$ is symmetric around zero,
ReLU sets half of the values to zero and keeps the other half, so
$\mathbb{E}[h_j^2] = \sigma_f^2/2$. Together,
\begin{equation}
  \sigma_{f'}^2 = \sigma_\Omega^2 \sum_{j=1}^{D_h} \frac{\sigma_f^2}{2}
  = \frac{1}{2} D_h\, \sigma_\Omega^2\, \sigma_f^2 .
  \label{eq:var-propagation}
\end{equation}
\noindent
Each layer therefore multiplies the variance by the factor
$\tfrac12 D_h \sigma_\Omega^2$. The variance stays constant exactly when this
factor is one, that is, when
\begin{equation}
  \sigma_\Omega^2 = \frac{2}{D_h} .
  \label{eq:he-forward}
\end{equation}

\paragraph{He initialization.}
The backward pass leads to the same argument with the roles of the layers
swapped. The gradient with respect to the activations of layer $k$ is a sum over
the $D_{h'}$ units of layer $k+1$, so keeping the gradient variance constant
requires $\sigma_\Omega^2 = 2/D_{h'}$. When the two layers have different sizes,
the two conditions cannot both hold, and the usual compromise uses the mean of
$D_h$ and $D_{h'}$:
\begin{equation}
  \sigma_\Omega^2 = \frac{2}{(D_h + D_{h'})/2} = \frac{4}{D_h + D_{h'}} .
  \label{eq:he-init}
\end{equation}
\noindent
This rule is called \emph{He} or \emph{Kaiming initialization}. It assumes ReLU
activations; the factor $2$ compensates for the half of the signal that ReLU
removes. For convolutional layers, $D_h$ is the number of inputs of one output
unit, $C_i K^2$ for kernel size $K$ and $C_i$ input channels.

\paragraph{Worked example: fifty layers of width 100.}
Consider a network with $50$ hidden layers of $D_h = 100$ units each. By
Eq.~\eqref{eq:var-propagation} the variance is multiplied by
$\tfrac12 \cdot 100 \cdot \sigma_\Omega^2 = 50\,\sigma_\Omega^2$ per layer, so
after $50$ layers by $(50\,\sigma_\Omega^2)^{50}$.
He initialization gives $\sigma_\Omega^2 = 2/100 = 0.02$ and a factor of exactly
$1$. Choosing $\sigma_\Omega^2 = 0.1$ looks harmless, but gives a factor of $5$
per layer and $5^{50} \approx 10^{35}$ overall; $\sigma_\Omega^2 = 0.01$ gives
$0.5^{50} \approx 10^{-15}$. Figure~\ref{fig:init-variance} shows these growth
curves. The backward pass behaves in the same way, with the gradient variance
growing or shrinking from the output back to the input. With $32$-bit floating
point numbers (range about $10^{\pm38}$), the settings far from $0.02$ make
training impossible before it begins.

\begin{figure}[ht]
  \centering
  \begin{tikzpicture}[x=0.1cm, y=0.035cm]
    \draw[-{Latex[length=2mm]}] (0,-80) -- (0,100);
    \draw[-{Latex[length=2mm]}] (0,-80) -- (58,-80);
    \foreach \t in {-50,0,50} {
      \draw (-1,\t) -- (1,\t);
      \node[left, font=\small] at (-1,\t) {$10^{\t}$};
    }
    \foreach \t in {0,25,50} {
      \draw (\t,-81.5) -- (\t,-78.5);
      \node[below=2pt, font=\small] at (\t,-82) {$\t$};
    }
    \node[font=\small] at (29,-106) {layer $k$};
    \node[rotate=90, font=\small] at (-15,10) {variance of $h_k$};
    \draw[very thick, red!70!black] (0,0) -- (50,84.9);
    \draw[very thick, orange!80!black] (0,0) -- (50,34.9);
    \draw[very thick, gray] (0,0) -- (50,0);
    \draw[very thick, teal] (0,0) -- (50,-15.1);
    \draw[very thick, blue!70!black] (0,0) -- (50,-65.1);
    \node[right, font=\small] at (50,84.9) {$\sigma_\Omega^2 = 1$};
    \node[right, font=\small] at (50,34.9) {$0.1$};
    \node[right, font=\small] at (50,2) {$0.02$ (He)};
    \node[right, font=\small] at (50,-17) {$0.01$};
    \node[right, font=\small] at (50,-65.1) {$0.001$};
  \end{tikzpicture}
  \caption{Variance of the activations in a network with $50$ layers of $100$
  ReLU units, on a logarithmic scale, for five weight variances. Only He
  initialization, $\sigma_\Omega^2 = 2/D_h = 0.02$, keeps the variance constant;
  all other choices change it by many orders of magnitude.}
  \label{fig:init-variance}
\end{figure}

\noindent
Initialization only controls the network at the start of training; once the
weights change, the variances drift again. Pre-trained weights, discussed for
transfer learning in Section~\ref{sec:cnn}, are a different kind of
initialization: they start from weights that already encode useful features
instead of from random numbers with the right scale.

\subsection{Why Deeper Plain Networks Fail}

A network processes its input sequentially. Each layer receives the output of
the previous layer, so a network with four layers is a chain of nested
functions,
\begin{equation}
  \mathbf{h}_1 = \mathbf{f}_1[\mathbf{x}, \boldsymbol{\phi}_1],\quad
  \mathbf{h}_2 = \mathbf{f}_2[\mathbf{h}_1, \boldsymbol{\phi}_2],\quad
  \mathbf{h}_3 = \mathbf{f}_3[\mathbf{h}_2, \boldsymbol{\phi}_3],\quad
  \mathbf{y} = \mathbf{f}_4[\mathbf{h}_3, \boldsymbol{\phi}_4],
\end{equation}
or in one line
$\mathbf{y} = \mathbf{f}_4\bigl[\mathbf{f}_3\bigl[\mathbf{f}_2\bigl[\mathbf{f}_1[\mathbf{x}]\bigr]\bigr]\bigr]$.
Here $\mathbf{f}_k$ may be a single layer or a group of layers, and
$\boldsymbol{\phi}_k$ are its parameters.

\paragraph{Degradation: training error rises with depth.}
Even with good initialization, adding layers helps only up to a point. On the
CIFAR-10 dataset, a plain convolutional network with $56$ layers has a higher
error than one with $20$ layers, and this holds not only for the test error but
already for the \emph{training} error. This is not overfitting. An overfitting
model fits the training set too well; here the deeper model fits it worse,
although it could in principle copy the $20$-layer network and set its extra
layers to the identity. The deeper network is harder to \emph{optimize}.

\paragraph{Shattered gradients.}
A likely reason is that the gradients of deep networks are \emph{shattered}.
Gradient descent takes finite steps and relies on the gradient being roughly
constant over the length of one step. In a shallow network, the gradient with
respect to an early-layer parameter changes smoothly as the parameter changes.
In a deep network it changes quickly and unpredictably, even for tiny parameter
changes, and looks almost like noise; a step along the gradient then no longer
reliably reduces the loss. The chain rule shows where this comes from. The
derivative of the output with respect to the first layer is
\begin{equation}
  \frac{\partial \mathbf{y}}{\partial \mathbf{f}_1}
  = \frac{\partial \mathbf{f}_4}{\partial \mathbf{f}_3}\,
    \frac{\partial \mathbf{f}_3}{\partial \mathbf{f}_2}\,
    \frac{\partial \mathbf{f}_2}{\partial \mathbf{f}_1} .
  \label{eq:plain-chain}
\end{equation}
\noindent
Every factor is evaluated at activations that are computed from $\mathbf{f}_1$.
When the parameters of the first layer change, all factors in the product change
at once, and the product of many changing matrices varies in an increasingly
complex way the deeper the network is. Changes in early layers therefore modify
the output in a way that becomes harder and harder to predict.

\subsection{Residual Connections}

Residual connections, also called skip connections, fix this by changing what
each block computes. Instead of replacing its input, each block adds a
correction to it: the input of the block is added to the block's output.
\begin{equation}
  \begin{aligned}
    \mathbf{h}_1 &= \mathbf{x} + \mathbf{f}_1[\mathbf{x}, \boldsymbol{\phi}_1], &
    \mathbf{h}_2 &= \mathbf{h}_1 + \mathbf{f}_2[\mathbf{h}_1, \boldsymbol{\phi}_2], \\
    \mathbf{h}_3 &= \mathbf{h}_2 + \mathbf{f}_3[\mathbf{h}_2, \boldsymbol{\phi}_3], &
    \mathbf{y} &= \mathbf{h}_3 + \mathbf{f}_4[\mathbf{h}_3, \boldsymbol{\phi}_4].
  \end{aligned}
  \label{eq:residual}
\end{equation}
\noindent
The block $\mathbf{f}_k$ now only has to learn the \emph{residual}, the
difference between the desired output and the input. A block that should do
nothing only needs weights close to zero, which is easy to reach, whereas a
plain block would have to learn the identity map exactly. Input and output of a
residual block must have the same size so that they can be added.
Figure~\ref{fig:residual-chain} shows the resulting network.

\begin{figure}[ht]
  \centering
  \begin{tikzpicture}[
      blk/.style={draw, rounded corners, fill=blue!8, minimum width=0.9cm,
        minimum height=0.7cm},
      add/.style={draw, circle, inner sep=0pt, minimum size=0.42cm},
      arrow/.style={-{Latex[length=2mm]}, thick}
    ]
    \node (x) at (0,0) {$\mathbf{x}$};
    \foreach \i in {1,...,4} {
      \node[blk] (f\i) at ({2.6*\i-1.0},0) {$\mathbf{f}_\i$};
      \node[add] (p\i) at ({2.6*\i+0.1},0) {$+$};
      \draw[arrow] (f\i) -- (p\i);
      \coordinate (s\i) at ({2.6*\i-1.95},0);
      \draw[arrow] (s\i) -- ++(0,0.8) -| (p\i.north);
      \fill (s\i) circle (1.5pt);
    }
    \node (y) at (11.5,0) {$\mathbf{y}$};
    \draw[arrow] (x) -- (f1);
    \draw[arrow] (p1) -- (f2);
    \draw[arrow] (p2) -- (f3);
    \draw[arrow] (p3) -- (f4);
    \draw[arrow] (p4) -- (y);
    \node[below=0.2cm, font=\small] at ($(p1)!0.5!(f2)$) {$\mathbf{h}_1$};
    \node[below=0.2cm, font=\small] at ($(p2)!0.5!(f3)$) {$\mathbf{h}_2$};
    \node[below=0.2cm, font=\small] at ($(p3)!0.5!(f4)$) {$\mathbf{h}_3$};
    \node[above=0.85cm, font=\small] at ($(s1)!0.5!(p1)$) {identity};
  \end{tikzpicture}
  \caption{A residual network with four blocks. Each skip connection carries the
  block input unchanged around the block and adds it to the block output, so
  every block computes a correction to its input.}
  \label{fig:residual-chain}
\end{figure}

\paragraph{A residual network is a sum of many paths.}
Substituting the equations~\eqref{eq:residual} into each other unravels the
network:
\begin{equation}
  \begin{aligned}
  \mathbf{y} = \mathbf{x}
    &+ \mathbf{f}_1[\mathbf{x}]
     + \mathbf{f}_2\bigl[\mathbf{x} + \mathbf{f}_1[\mathbf{x}]\bigr]
     + \mathbf{f}_3\Bigl[\mathbf{x} + \mathbf{f}_1[\mathbf{x}]
       + \mathbf{f}_2\bigl[\mathbf{x} + \mathbf{f}_1[\mathbf{x}]\bigr]\Bigr] \\
    &+ \mathbf{f}_4\Bigl[\mathbf{x} + \mathbf{f}_1[\mathbf{x}]
       + \mathbf{f}_2\bigl[\mathbf{x} + \mathbf{f}_1[\mathbf{x}]\bigr]
       + \mathbf{f}_3\bigl[\ldots\bigr]\Bigr] .
  \end{aligned}
\end{equation}
\noindent
The output is the input plus the contributions of networks of every depth from
one to four. Each block can be passed through or skipped, so $K$ blocks give
$2^K$ paths from input to output: for $K = 4$, one path of length $0$ (the pure
identity), four of length $1$, six of length $2$, four of length $3$, and one of
length $4$. A residual network behaves like an ensemble of shallower networks,
and most of its paths are much shorter than the full depth.

\paragraph{Gradients through the skip connection.}
The same structure appears in the backward pass. The Jacobian of one residual
block is
\begin{equation}
  \frac{\partial \mathbf{h}_k}{\partial \mathbf{h}_{k-1}}
  = \mathbf{I} + \frac{\partial \mathbf{f}_k}{\partial \mathbf{h}_{k-1}} ,
\end{equation}
so the derivative of the output with respect to the input is the product
$\prod_k \bigl(\mathbf{I} + \partial \mathbf{f}_k/\partial \mathbf{h}_{k-1}\bigr)$.
Multiplying it out gives again $2^K$ terms, one per path, and the first term is
the identity. Unlike Eq.~\eqref{eq:plain-chain}, the gradient is no longer a
single long product: the short paths contribute terms with few factors, which
change slowly and predictably when the early layers change. In the language of
computational graphs, the skip connection ends in an addition node, and an
addition node distributes the upstream gradient unchanged to both of its inputs
(Table~\ref{tab:gate-rules}). It does not send its inputs backward. The gradient
of the loss therefore reaches every block directly through the chain of skip
connections, however deep the network is.

\paragraph{Order of operations inside a block.}
Where the nonlinearity sits in a residual block matters:
\begin{itemize}
  \item If a block ends with ReLU, as in linear layer followed by ReLU, then
        $\mathbf{f}_k \ge 0$ and each block can only \emph{increase} the values
        of its input. The ReLU is therefore placed at the beginning of the
        block: ReLU, then linear layer.
  \item If every block starts with ReLU, a block does nothing for negative
        inputs. The network therefore begins with a plain linear transformation
        before the first residual block, so that the signal entering the
        residual chain is a learned mixture rather than the raw input.
  \item A block usually contains several layers, for example ReLU, linear,
        ReLU, linear, before the skip connection adds the input back.
\end{itemize}

\begin{table}[!ht]
  \centering
  \begin{tabular}{@{}lccccl@{}}
    \toprule
    Variant & $\mathbf{x}$ & $\mathbf{h}_1$ & $\mathbf{h}_2$ & $\mathbf{h}_3$ &
      after $K$ blocks \\
    \midrule
    plain residual blocks & $1$ & $2$ & $4$ & $8$ & $2^K$ \\
    output scaled by $1/\sqrt{2}$ & $1$ & $1$ & $1$ & $1$ & $1$ \\
    batch normalization at block start & $1$ & $2$ & $3$ & $4$ & $K+1$ \\
    \bottomrule
  \end{tabular}
  \caption{Variance of the signal after each residual block when input and
  block output both have unit variance. Plain addition doubles the variance per
  block; scaling by $1/\sqrt{2}$ keeps it constant; batch normalization inside
  the residual branch lets it grow only linearly.}
  \label{tab:resnet-variance}
\end{table}

\paragraph{Variance in residual networks.}
Residual connections solve the shattered gradients, but they create a new
problem for the forward pass. Suppose the input of a block has variance $1$ and
He initialization makes the block output $\mathbf{f}_k$ also have variance $1$.
If the two are uncorrelated, their sum has variance $1 + 1 = 2$, so the variance
\emph{doubles} with every block (Table~\ref{tab:resnet-variance}), and the
network again suffers from unstable forward propagation and exploding
gradients. One remedy combines He initialization with a multiplication of each
block's combined output by $1/\sqrt{2}$, which divides the variance by two and
restores it to $1$. More common is batch normalization, discussed next. With
these measures, residual connections roughly double the depth that can be
trained before performance degrades; together with normalization, networks with
hundreds of layers become routine.

\subsection{Batch Normalization}

Batch normalization (BatchNorm, BN) attacks the variance problem directly.
Instead of hoping that the initialization keeps activations in a good range, it
measures their statistics during training and rescales them. For every hidden
unit, the activations are standardized across the training examples of the
current mini-batch $\mathcal{B}$. First the empirical mean and standard
deviation of the unit's activation $h$ over the batch are computed,
\begin{equation}
  m_h = \frac{1}{|\mathcal{B}|} \sum_{i \in \mathcal{B}} h_i , \qquad
  s_h = \sqrt{\frac{1}{|\mathcal{B}|} \sum_{i \in \mathcal{B}} (h_i - m_h)^2} ,
\end{equation}
then each activation is standardized to mean $0$ and variance $1$,
\begin{equation}
  h_i \leftarrow \frac{h_i - m_h}{s_h + \epsilon}
  \qquad \forall i \in \mathcal{B} ,
\end{equation}
where the small constant $\epsilon$ prevents division by zero. Finally a learned
scale $\gamma$ and offset $\delta$ give the network back the freedom to choose
the mean and standard deviation itself:
\begin{equation}
  h_i \leftarrow \gamma h_i + \delta \qquad \forall i \in \mathcal{B} .
  \label{eq:bn-affine}
\end{equation}
\noindent
The same layer is often written as
$\hat{x} = (x - \mu_\mathcal{B})/\sqrt{\sigma_\mathcal{B}^2 + \epsilon}$,
$y = \gamma \hat{x} + \beta$, with batch mean $\mu_\mathcal{B}$, batch variance
$\sigma_\mathcal{B}^2$, and offset $\beta$ instead of $\delta$; the two forms
differ only in where $\epsilon$ enters. Without $\gamma$ and $\delta$, the layer
would force every unit to zero mean and unit variance forever. With them, the
network can learn any scale, but the scale is now an explicit parameter of the
layer instead of an accident of the weights of all preceding layers.

\paragraph{Where the statistics are computed.}
Batch normalization is applied independently to each hidden unit, and each unit
gets its own $\gamma$ and $\delta$. A fully connected network with $K$ layers of
$D$ hidden units therefore learns $KD$ offsets and $KD$ scales. In a
convolutional network, the statistics are computed not only across the batch but
also across all spatial positions of a channel, in keeping with weight sharing:
every position of a feature map is treated in the same way. A convolutional
network with $K$ layers of $C$ channels thus learns $KC$ offsets and $KC$
scales. For example, a layer with $64$ channels, a batch of $32$ images, and
$56\times56$ feature maps computes each of its $64$ means over
$32\cdot56\cdot56 = 100{,}352$ values and adds $128$ learned parameters.

\paragraph{Training mode and evaluation mode.}
At test time there may be only a single image, and the prediction for one image
should not depend on which other images happen to be in the same batch. Batch
normalization therefore behaves differently in the two phases. During training
it uses the statistics of the current batch and additionally keeps running
averages of the means and variances. During evaluation it uses these stored
averages as fixed constants, so the layer becomes a fixed affine map. Forgetting
to switch a model into evaluation mode for validation and testing is a classic
error: the outputs then depend on the batch composition.

\paragraph{Invariance to the scale of the weights.}
If all weights and biases that feed a batch-normalized unit are multiplied by
some $\alpha > 0$, the mean and standard deviation over the batch are multiplied
by $\alpha$ as well, and the standardized value does not change. Batch
normalization thus makes the network invariant to rescaling weights and biases,
because it compensates for it. This removes exactly the sensitivity to scale
that made initialization so critical.

\paragraph{Other normalizations.}
Batch normalization depends on the batch, which is a disadvantage for small
batches and for sequence models. Other normalizations use the same
standardize-then-scale recipe but compute the statistics over different axes of
the activation tensor with batch size $N$, channels $C$, and spatial size
$H\times W$ (Table~\ref{tab:norm-variants}). Layer normalization, instance
normalization, and group normalization compute their statistics within a single
example and therefore behave identically in training and evaluation.

\begin{table}[ht]
  \centering
  \begin{tabular}{@{}p{2.8cm}p{5.6cm}p{5.6cm}@{}}
    \toprule
    Method & Statistics computed over & One mean per \\
    \midrule
    Batch norm & $N$, $H$, $W$ & channel (shared by all examples) \\
    Layer norm & $C$, $H$, $W$ & example \\
    Instance norm & $H$, $W$ & example and channel \\
    Group norm & $H$, $W$, and a group of channels & example and channel group \\
    \bottomrule
  \end{tabular}
  \caption{Normalization layers differ only in the axes of the $N\times C\times
  H\times W$ activation tensor over which mean and variance are computed. Only
  batch normalization mixes different examples.}
  \label{tab:norm-variants}
\end{table}

\paragraph{Costs and benefits.}
Batch normalization has three main effects:
\begin{itemize}
  \item \emph{Stable forward propagation.} With the initialization
        $\gamma = 1$, $\delta = 0$, every normalized activation has unit
        variance. In a network without residual connections this keeps the
        variance constant through all layers. In a residual network with batch
        normalization at the start of each residual branch, each block adds
        variance $1$ to the signal, so the variance grows as $1, 2, 3, 4, \ldots$
        (Table~\ref{tab:resnet-variance}). Later blocks then contribute a
        smaller share to the sum, and the network is effectively less deep at
        the start of training. As training proceeds, $\gamma$ can grow; the
        network thereby controls its own effective depth.
  \item \emph{Higher learning rates.} Batch normalization makes the loss
        surface and its gradients change more smoothly, which allows larger
        learning rates and faster training.
  \item \emph{Regularization.} The normalization depends on the other examples
        in the batch, so every activation receives a small amount of random
        noise during training. Such noise can improve generalization, similar
        to dropout.
\end{itemize}
\noindent
The costs are the dependence on the batch: small batches give noisy
statistics, and training and evaluation behave differently.

\subsection{Common Residual Architectures}

Residual connections, He initialization, and batch normalization are the building
blocks of nearly all deep image networks. Four architectures show how they are
combined: ResNet for classification, DenseNet as a variant with concatenation,
U-Net for segmentation, and stacked hourglass networks for pose estimation.

\paragraph{ResNet blocks.}
The standard ResNet block for images applies batch normalization, then ReLU,
then a $3\times3$ convolution, and repeats this sequence once before the skip
connection adds the input. This follows the order of operations above:
normalization and nonlinearity come first, and the block ends with a linear
operation. Deep ResNets use the \emph{bottleneck block} instead
(Figure~\ref{fig:bottleneck}). A $1\times1$ convolution first reduces the number
of channels by a factor of four, the expensive $3\times3$ convolution then works
on the reduced representation, and a second $1\times1$ convolution increases the
number of channels again so that the input can be added.

\begin{figure}[ht]
  \centering
  \begin{tikzpicture}[
      blk/.style={draw, rounded corners, fill=blue!8, align=center,
        minimum height=0.95cm, minimum width=2.4cm, font=\small},
      add/.style={draw, circle, inner sep=0pt, minimum size=0.42cm},
      arrow/.style={-{Latex[length=2mm]}, thick}
    ]
    \node[blk] (b1) at (1.9,0) {BN, ReLU,\\conv $1\times1$};
    \node[blk] (b2) at (5.3,0) {BN, ReLU,\\conv $3\times3$};
    \node[blk] (b3) at (8.7,0) {BN, ReLU,\\conv $1\times1$};
    \node[add] (p) at (11.0,0) {$+$};
    \coordinate (in) at (0,0);
    \coordinate (s) at (0.35,0);
    \draw[arrow] (in) -- (b1);
    \draw[arrow] (b1) -- (b2);
    \draw[arrow] (b2) -- (b3);
    \draw[arrow] (b3) -- (p);
    \draw[arrow] (p) -- (12.0,0);
    \draw[arrow] (s) -- ++(0,1.0) -| (p.north);
    \fill (s) circle (1.5pt);
    \node[below=0.15cm, font=\small] at (0.35,0) {$256$};
    \node[below=0.15cm, font=\small] at ($(b1.east)!0.5!(b2.west)$) {$64$};
    \node[below=0.15cm, font=\small] at ($(b2.east)!0.5!(b3.west)$) {$64$};
    \node[below=0.15cm, font=\small] at ($(b3.east)!0.5!(p.west)$) {$256$};
    \node[below=0.55cm of b1, font=\small] {reduce channels};
    \node[below=0.55cm of b3, font=\small] {increase channels};
  \end{tikzpicture}
  \caption{Bottleneck ResNet block. Numbers below the arrows are channel
  counts. The two $1\times1$ convolutions reduce the channels by a factor of
  four and restore them, so the $3\times3$ convolution runs on only $64$
  channels.}
  \label{fig:bottleneck}
\end{figure}

\noindent
The saving is large. With $256$ input and output channels and ignoring biases, a
standard block with two $3\times3$ convolutions needs
$2 \cdot 9 \cdot 256^2 \approx 1.18$ million weights. The bottleneck block needs
$256\cdot64 + 9\cdot64^2 + 64\cdot256 = 16{,}384 + 36{,}864 + 16{,}384
= 69{,}632$ weights, about seventeen times fewer, and it contains three
nonlinear layers instead of two.

\paragraph{ResNet-200.}
The ResNet-200 network for ImageNet starts with a $7\times7$ convolution with
stride $2$ and a max pooling step, reducing the $224\times224\times3$ image to
$56\times56$ with $64$ channels. Four stages of bottleneck blocks follow, with
$56\times56\times256$, $28\times28\times512$, $14\times14\times1024$, and
$7\times7\times2048$ outputs; inside the blocks the channels are reduced to $64$,
$128$, $256$, and $512$. Between the stages the resolution is halved, and the
skip connection uses a strided $1\times1$ convolution so that input and output
shapes match again. The two middle stages contain $24$ and $36$ blocks; with
three blocks each in the first and last stage this gives $66$ blocks of three
convolutions, plus the first convolution and the final layer, $200$ layers in
total. Global average pooling, a fully connected layer, and a softmax over the
$1000$ classes produce the output. ResNet-200 reached a top-5 error of $4.8\%$ on
ImageNet, the first network to exceed human performance on this benchmark, and
continued the trend of Table~\ref{tab:alexnet-vgg} from $16.4\%$ (AlexNet)
and $6.8\%$ (VGG).

\paragraph{DenseNet.}
Instead of adding the input activations to the output, they can also be
\emph{concatenated} with it. In a DenseNet, the input of every layer consists of
the concatenated outputs of \emph{all} previous layers. Early layers then
contribute directly to the output, which, as with residual connections, makes
the loss surface behave reasonably. Concatenation lets the number of channels
grow with every layer, so a $1\times1$ convolution reduces it from time to
time. Concatenation across a downsampling step is impossible, because the
representations have different spatial sizes. The chain of concatenations is
therefore broken at each downsampling step, and the smaller representation
starts a new chain. With the same number of parameters, DenseNet can perform
better than ResNet.

\paragraph{U-Net.}
The encoder--decoder networks for semantic segmentation in Section~\ref{sec:cnn}
must squeeze the whole image through a low-resolution bottleneck. The decoder
then has to recover the exact position of every boundary from coarse features,
as if it had to ``remember'' the high-resolution detail. The U-Net adds skip
connections from each encoder level to the decoder level of the same
resolution (Figure~\ref{fig:unet}). The decoder receives the fine details
directly and only has to combine them with the semantic information from the
bottleneck. The U-Net uses valid $3\times3$ convolutions, so the encoder feature
maps are slightly larger than the corresponding decoder maps; they are
\emph{cropped} to the same size and then \emph{concatenated}, not added.
Upsampling uses transposed $2\times2$ convolutions, and a final $1\times1$
convolution produces one score per class and pixel. The input of
$572\times572$ pixels therefore gives a label map of $388\times388$ pixels.

\begin{figure}[ht]
  \centering
  \begin{tikzpicture}[
      box/.style={draw, fill=blue!10, minimum width=0.35cm},
      arrow/.style={-{Latex[length=2mm]}, thick},
      skip/.style={-{Latex[length=2mm]}, thick, orange!80!black}
    ]
    \node[box, minimum height=0.9cm] (e1) at (0,0) {};
    \node[box, minimum height=0.75cm] (e2) at (1.2,-1) {};
    \node[box, minimum height=0.6cm] (e3) at (2.4,-2) {};
    \node[box, minimum height=0.45cm] (e4) at (3.6,-3) {};
    \node[box, fill=orange!25, minimum height=0.3cm] (b) at (4.8,-4) {};
    \node[box, minimum height=0.45cm] (d4) at (6.0,-3) {};
    \node[box, minimum height=0.6cm] (d3) at (7.2,-2) {};
    \node[box, minimum height=0.75cm] (d2) at (8.4,-1) {};
    \node[box, minimum height=0.9cm] (d1) at (9.6,0) {};
    \foreach \a/\b in {e1/e2, e2/e3, e3/e4, e4/b} {
      \draw[arrow] (\a.south) -- (\b.north west);
    }
    \foreach \a/\b in {b/d4, d4/d3, d3/d2, d2/d1} {
      \draw[arrow] (\a.north east) -- (\b.south);
    }
    \foreach \i in {1,...,4} {
      \draw[skip] (e\i.east) -- (d\i.west);
    }
    \node[above=0.05cm, font=\small] at ($(e1.east)!0.5!(d1.west)$)
      {crop and concatenate};
    \node[left=0.1cm of e1, font=\small] {$568^2\times64$};
    \node[left=0.1cm of e2, font=\small] {$280^2\times128$};
    \node[left=0.1cm of e3, font=\small] {$136^2\times256$};
    \node[left=0.1cm of e4, font=\small] {$64^2\times512$};
    \node[below=0.1cm of b, font=\small] {$28^2\times1024$};
    \node[right=0.1cm of d4, font=\small] {$52^2\times512$};
    \node[right=0.1cm of d3, font=\small] {$100^2\times256$};
    \node[right=0.1cm of d2, font=\small] {$196^2\times128$};
    \node[right=0.1cm of d1, font=\small] {$388^2\times64$};
  \end{tikzpicture}
  \caption{U-Net. The encoder (left) halves the resolution with max pooling, the
  decoder (right) doubles it with transposed convolutions. Horizontal skip
  connections copy each encoder map, crop it, and concatenate it with the
  decoder map of the same level. Labels give spatial size and channels.}
  \label{fig:unet}
\end{figure}

\noindent
The U-Net was designed for segmenting medical images. A typical example is the
classification of voxels in three-dimensional volumes of mouse cortex recorded
with a scanning electron microscope, where a single U-Net already gives a good
labeling and an ensemble of five U-Nets a better one. Because its output has the
size of its input and preserves fine detail, the U-Net is also the standard
backbone for other image-to-image tasks, including the denoising networks of
diffusion models.

\paragraph{Stacked hourglass networks.}
A stacked hourglass network repeats several U-Net-like encoder--decoder blocks
one after the other. Each block alternates between looking at the image at a
local and at a global level, and the next block refines the result of the
previous one. Inside the blocks, skip connections \emph{add} the results rather
than concatenating them. The network was developed for human pose estimation.
It produces one \emph{heatmap} per body joint, whose value at each pixel
indicates how likely the joint is to lie there, and the estimated joint position
is the maximum of its heatmap. Training compares the predicted heatmaps with
target heatmaps, so pose estimation is formulated as a \emph{regression}
problem, in contrast to the classification problems considered so far.

\paragraph{Why residual networks work so well.}
Visualizations of the loss as a function of the parameters make the difference
visible. For a deep plain network, the loss surface is rugged and full of
local structure, matching the shattered gradients. With residual connections,
the loss surface around the solution is much smoother and nearly convex, so
gradient descent with a finite step size makes steady progress.
Table~\ref{tab:residual-archs} summarizes the four architectures.

\begin{table}[ht]
  \centering
  \begin{tabular}{@{}p{2.8cm}p{3.0cm}p{2.8cm}p{6.0cm}@{}}
    \toprule
    Architecture & Skip operation & Typical task & Key idea \\
    \midrule
    ResNet & addition & classification & bottleneck blocks with BN and ReLU;
      hundreds of layers \\
    DenseNet & concatenation & classification & every layer sees all previous
      outputs; $1\times1$ convolutions limit channels \\
    U-Net & crop and concatenate & segmentation & skips from encoder to decoder
      preserve high-resolution detail \\
    Stacked hourglass & addition & pose estimation & repeated U-Net-like blocks;
      one heatmap per joint (regression) \\
    \bottomrule
  \end{tabular}
  \caption{Four common residual architectures. They differ in how the skip
  connection combines its signals and in the task the network is built for.}
  \label{tab:residual-archs}
\end{table}
